\section{Periodic Construction and Floquet Reduction}

We now place the order-32 pattern in an arbitrary number of eight-site cells.
Let

\begin{lstlisting}[language={}]
tau_*=(+,+,-,+,-,-,+,-)                                         (4.1)
\end{lstlisting}

and impose either Hamilton holonomy \(\alpha \in {+1,-1}\) on a graph of order
\(n=8L\).

\subsection{The eight-site fiber}

Write \(i=8m+r\), \(0\le r<8\), and put \(x_{8m+r}=z^m v_r\). Substitution in (2.1)
gives

\begin{lstlisting}[language={}]
H(z)=
[ 0  1  1  0  0  0   z^-1   z^-1 ]
[ 1  0  1  1  0  0    0    -z^-1 ]
[ 1  1  0  1 -1  0    0      0   ]
[ 0  1  1  0  1  1    0      0   ]
[ 0  0 -1  1  0  1   -1      0   ]
[ 0  0  0  1  1  0    1     -1   ]
[ z  0  0  0 -1  1    0      1   ]
[ z -z  0  0  0 -1    1      0   ].                              (4.2)
\end{lstlisting}

On \(|z|=1\), replacing \(z\) by its complex conjugate \(z^{-1}\) transposes the
matrix, so (4.2) is Hermitian.

\textbf{Proposition 4.1 (finite Floquet decomposition).} For \(n=8L\),

\begin{lstlisting}[language={}]
A_(8L,alpha) ~= direct_sum_(z^L=alpha) H(z).                     (4.3)
\end{lstlisting}

\textbf{Proof.} Let \(S\) be the unitary shift on the \(L\) cell coordinates with
twisted boundary \(u_{m+L}=\alpha u_m\). Its normalized eigenvectors are
\(L^{-1/2}(1,z,...,z^{L-1})\) with \(z^L=\alpha\). Distinct roots are orthogonal by
the geometric-sum identity. The signed operator is a finite sum of residue
matrices tensored with powers of \(S\), so each eigenspace is invariant and its
restriction is (4.2). There are \(L\) such eight-dimensional eigenspaces, and
their dimensions sum to \(8L\). \(\square\)

Thus the finite problem uses the discrete set \(z^L=\alpha\), while the infinite
periodic problem uses every \(|z|=1\).

\subsection{Exact determinant}

Let \(x\) be the spectral parameter. Fraction-free elimination of
\(xI-H(z)\) gives

\begin{lstlisting}[language={}]
det(xI-H(z))
 =x^8-16x^6+(80-2(z+z^-1))x^4
   +(-128+16(z+z^-1))x^2
   +(z^2+z^-2)-13(z+z^-1)+40.                                  (4.4)
\end{lstlisting}

Set

\begin{lstlisting}[language={}]
y=x^2,       c=z+z^-1.
\end{lstlisting}

Since \(z^2+z^{-2}=c^2-2\), equation (4.4) becomes

\begin{lstlisting}[language={}]
P(y,c)=y^4-16y^3+(80-2c)y^2+(-128+16c)y+c^2-13c+38.             (4.5)
\end{lstlisting}

For \(|z|=1\), \(c=2\operatorname{cos}(\theta)\) lies in \([-2,2]\). Hermiticity implies that every
eigenvalue \(\lambda\) is real, and (4.5) shows that \(\lambda^2\) is a nonnegative
root of \(P(y,c)\).

\subsection{A uniform rational bound}

Put

\begin{lstlisting}[language={}]
B=1561/200,       u=y-B,       t=2-c.
\end{lstlisting}

Direct substitution in (4.5) gives

\begin{lstlisting}[language={}]
P(B+u,2-t)
 =u^4+(761/50)u^3+2u^2t+(1337363/20000)u^2
  +(761/50)ut+(136311081/2000000)u
  +t^2+(119121/20000)t+84332641/1600000000.                     (4.6)
\end{lstlisting}

Every coefficient in (4.6) is nonnegative and the constant term is positive.
Therefore \(P(y,c)>0\) whenever \(y\ge B\) and \(c\le 2\). In particular, no squared
eigenvalue of any unit-circle fiber reaches \(B\); hence

\begin{lstlisting}[language={}]
rho(H(z))^2<B                                                        (4.7)
\end{lstlisting}

for every \(|z|=1\). Proposition 4.1 now yields

\begin{lstlisting}[language={}]
rho(A_(8L,alpha))^2<1561/200                                      (4.8)
\end{lstlisting}

for both holonomies and every \(L\ge 1\).

\subsection{Comparison with the conjectured threshold}

At \(n=32\), the Sturm isolation (3.4)-(3.5) gives

\begin{lstlisting}[language={}]
rho_-(32)^2>7809/1000>1561/200.                                  (4.9)
\end{lstlisting}

For \(n\ge 32\), both angles \(2\pi/n\) and \(4\pi/n\) lie in \((0,\pi)\) and decrease as
\(n\) increases. Since cosine decreases with its argument on \((0,\pi)\), equation
(2.12) shows that \(\rho_-(n)^2\) increases with \(n\). Thus

\begin{lstlisting}[language={}]
rho_-(n)^2>=rho_-(32)^2>1561/200.                                (4.10)
\end{lstlisting}

Combining (4.8) and (4.10) proves Theorem B for \(n=8L\), \(L\ge 4\). \(\square\)

The restriction to multiples of eight comes from the construction, not from
the threshold comparison. No conclusion is drawn here for other even orders.
